% Options for packages loaded elsewhere
\PassOptionsToPackage{unicode}{hyperref}
\PassOptionsToPackage{hyphens}{url}
\documentclass[
]{article}
\usepackage{xcolor}
\usepackage{amsmath,amssymb}
\setcounter{secnumdepth}{-\maxdimen} % remove section numbering
\usepackage{iftex}
\ifPDFTeX
  \usepackage[T1]{fontenc}
  \usepackage[utf8]{inputenc}
  \usepackage{textcomp} % provide euro and other symbols
  % Feature-completeness table symbols (✓ / ◐ are not in pdflatex fonts)
  \DeclareUnicodeCharacter{2713}{$\checkmark$}
  \DeclareUnicodeCharacter{25D0}{$\circ$}
\else % if luatex or xetex
  \usepackage{unicode-math} % this also loads fontspec
  \defaultfontfeatures{Scale=MatchLowercase}
  \defaultfontfeatures[\rmfamily]{Ligatures=TeX,Scale=1}
\fi
\usepackage{lmodern}
\ifPDFTeX\else
  % xetex/luatex font selection
\fi
% Use upquote if available, for straight quotes in verbatim environments
\IfFileExists{upquote.sty}{\usepackage{upquote}}{}
\IfFileExists{microtype.sty}{% use microtype if available
  \usepackage[]{microtype}
  \UseMicrotypeSet[protrusion]{basicmath} % disable protrusion for tt fonts
}{}
\makeatletter
\@ifundefined{KOMAClassName}{% if non-KOMA class
  \IfFileExists{parskip.sty}{%
    \usepackage{parskip}
  }{% else
    \setlength{\parindent}{0pt}
    \setlength{\parskip}{6pt plus 2pt minus 1pt}}
}{% if KOMA class
  \KOMAoptions{parskip=half}}
\makeatother
\usepackage{longtable,booktabs,array}
\usepackage{caption}
\captionsetup[table]{skip=6pt}
\newcounter{none} % for unnumbered tables
\usepackage{calc} % for calculating minipage widths
% Correct order of tables after \paragraph or \subparagraph
\usepackage{etoolbox}
\makeatletter
\patchcmd\longtable{\par}{\if@noskipsec\mbox{}\fi\par}{}{}
\makeatother
% Allow footnotes in longtable head/foot
\IfFileExists{footnotehyper.sty}{\usepackage{footnotehyper}}{\usepackage{footnote}}
\makesavenoteenv{longtable}
\setlength{\emergencystretch}{3em} % prevent overfull lines
\providecommand{\tightlist}{%
  \setlength{\itemsep}{0pt}\setlength{\parskip}{0pt}}
\usepackage[]{natbib}
\bibliographystyle{plainnat}
\usepackage{bookmark}
\IfFileExists{xurl.sty}{\usepackage{xurl}}{} % add URL line breaks if available
\urlstyle{same}
% fallback for those not using the hyperref driver hyperxmp:
\makeatletter
\@ifundefined{xmpquote}{\newcommand{\xmpquote}[1]{#1}}{}
\makeatother
\hypersetup{
  pdftitle={maxentcpp: A C++ reimplementation of Maximum Entropy species distribution modeling for R},
  hidelinks,
  pdfcreator={LaTeX via pandoc}}

\title{maxentcpp: A C++ reimplementation of Maximum Entropy species
distribution modeling for R}
\author{Ángel Luis Robles Fernández\thanks{Corresponding author. Department of Ecology and Evolutionary Biology, University of Kansas, Lawrence, KS, United States. ORCID: 0000-0002-4741-8012. E-mail: a.l.robles.fernandez@gmail.com}}
\date{11 August 2026}

\begin{document}
\maketitle

\section{Summary}\label{summary}

Maximum Entropy (Maxent) modeling is the \emph{de facto} standard for
presence-only species distribution modeling (SDM)
\citep{Phillips2006, Phillips2004, Elith2011}, with foundational papers
cited tens of thousands of times across ecology, conservation biology,
and biogeography. Yet the reference implementation remains a Java
desktop application \citep{Phillips2017}, and every species run in a
typical workflow --- e.g., 500 species \(\times\) 4 climatic scenarios
--- spawns a JVM and round-trips data through SWD files on disk.
\texttt{maxentcpp} is an R package that removes this friction with a
native C++17 reimplementation of the Maxent algorithm. It estimates the
geographic distribution of a species from presence-only occurrence
records and environmental covariates by identifying the probability
distribution of maximum entropy \citep{Jaynes1957} subject to
constraints derived from training data, in memory, with no Java runtime
and no file serialization.

The fitted model is a \emph{Gibbs distribution} over geographic space:

\[
P(\mathbf{x}) = \frac{1}{Z}\exp\!\left(\sum_{j=1}^{J} \lambda_j f_j(\mathbf{x})\right),
\]

where \(f_j\) are feature transformations of the environmental
variables, \(\lambda_j\) are learned weights, and
\(Z = \sum_{\mathbf{x}}\exp\!\bigl(\sum_j \lambda_j f_j(\mathbf{x})\bigr)\)
is the normalizing constant (partition function). The optimizer finds
the \(\lambda_j\) that maximize the regularized log-likelihood via
sequential coordinate ascent with \(\ell_1\) (lasso) penalties
\citep{Phillips2006, Elith2011}.

The original Maxent software is a Java desktop application
\citep{Phillips2017}. \texttt{maxentcpp} translates the core
continuous-predictor Maxent 3.4.4 fitting and projection algorithm into
C++17 with R bindings via Rcpp \citep{Eddelbuettel2011} and RcppEigen
\citep{Bates2013}, eliminating the Java Runtime Environment (JRE)
dependency entirely. The package supports all six feature types of Java
Maxent (linear, quadratic, product, hinge, threshold, and categorical
via \texttt{BinaryFeature}), output transformations in raw, logistic,
and complementary log-log (cloglog) scales, and a streaming raster
evaluation engine that processes environmental layers block-by-block
through the \texttt{terra} package \citep{Hijmans2024}, enabling
large-raster projections without loading entire raster stacks into
memory. Tools for diagnosing model performance include AUC evaluation,
permutation importance, response curves, jackknife variable-importance
analysis, replicate runs and stratified cross-validation, missing-data
handling, and Multivariate Environmental Similarity Surfaces (MESS)
\citep{Elith2010} for novelty detection.

\section{Statement of need}\label{statement-of-need}

Maxent is the \emph{de facto} standard for presence-only species
distribution modeling \citep{Merow2013}. However, the original Java
implementation presents a number of practical barriers for present-day
ecological research workflows:

\begin{enumerate}
\def\labelenumi{\arabic{enumi}.}
\item
  \textbf{Java dependency.} Java and \texttt{rJava} configuration can be
  a recurring source of installation and runtime failures, particularly
  across operating systems, Java versions, and managed computing
  environments. While \texttt{conda} environments and Docker containers
  can manage Java versioning, they add deployment complexity and are not
  standard practice in many ecology research groups.
\item
  \textbf{File-based I/O.} Data must be serialized to disk as SWD or
  ASCII grid files before the Java process can read them, and results
  must be parsed back from output files. There is no in-memory bridge
  between R data structures and the Java optimizer.
\item
  \textbf{Scalability.} The Java implementation is single-threaded and
  GUI-centric. In multi-species, multi-scenario workflows (e.g., 500
  species \(\times\) 4 climatic scenarios), file I/O overhead
  accumulates: each species run requires writing SWD files, invoking the
  JVM, and parsing output files. A native in-memory API removes this
  per-call overhead by keeping data and model state in R objects.
\item
  \textbf{Maintenance.} The Java Maxent codebase has received only minor
  updates since version 3.4.4, with limited active development
  \citep{Phillips2017}.
\end{enumerate}

\texttt{maxentcpp} addresses these barriers by providing a native C++/R
implementation. At present, the development version can be installed
from GitHub using
\texttt{remotes::install\_github("alrobles/maxentcpp")}; following CRAN
acceptance, the package will be installable with
\texttt{install.packages("maxentcpp")}. The package operates entirely in
memory through R objects and integrates seamlessly with the
\texttt{terra} spatial data ecosystem. The intended users are
ecologists, conservation biologists, and biogeographers who employ
Maxent in reproducible, script-based workflows --- from individual
species analyses to large-scale biodiversity assessments.

\section{Related work}\label{related-work}

\texttt{maxentcpp} fits within a broader modernization trend in
ecological and environmental-science software. Circuitscape.jl
\citep{Hall2021} illustrates how established ecological algorithms can
be reimplemented in a modern high-performance language to improve
scalability and parallel execution, while retaining continuity with a
widely used landscape-connectivity tool. In R, the spatial ecosystem has
similarly moved from older \texttt{sp}/\texttt{raster}-centered
workflows toward \texttt{sf} \citep{Pebesma2018} and \texttt{terra}
\citep{Hijmans2024}, and SDM packages such as \texttt{ENMeval}
\citep{Kass2021} and \texttt{biomod2} \citep{Thuiller2009} have followed
this transition --- with \texttt{ENMeval} additionally removing
Java/\texttt{rJava} from its default path in favor of \texttt{maxnet}.

Within Maxent workflows specifically, existing modernization efforts
address different parts of the problem. \textbf{dismo}
\citep{Hijmans2023} is the most widely used R interface to Java Maxent;
it wraps \texttt{maxent.jar} via \texttt{rJava}, inheriting the JDK
dependency and file I/O overhead, and is currently in maintenance mode
following the transition from \texttt{raster} to \texttt{terra}.
\textbf{rmaxent} \citep{Baumgartner2017} provides Java-free projection
of previously fitted Maxent models from \texttt{.lambdas} files, parses
feature weights and normalizers, and implements MESS-related
diagnostics, but does not replace Java for model fitting.
\textbf{maxnet} \citep{Phillips2024maxnet, Friedman2010} provides a
complete Java-free fitting implementation based on \texttt{glmnet}
coordinate descent --- preserving the same statistical model but
employing a different optimization algorithm that can produce
numerically different results in some settings. \textbf{ENMeval}
\citep{Kass2021} is a model tuning and evaluation framework that wraps
either \texttt{dismo::maxent()} or \texttt{maxnet::maxnet()}; it does
not implement the Maxent algorithm itself. \textbf{kuenm}
\citep{Cobos2019} is a calibration and evaluation toolkit that relies on
\texttt{dismo} or direct Java Maxent calls.

\texttt{maxentcpp} occupies a distinct position within this landscape by
porting the Java Maxent \texttt{density.Sequential} training optimizer
itself to C++17 and validating optimizer trajectories against the Java
implementation. To the best of current knowledge, \texttt{maxentcpp} is
the first R package to target the Java Maxent 3.4.4
\texttt{density.Sequential} optimizer through a native C++ port and to
publish per-iteration trajectory tests via the companion package
\texttt{maxentcppCompTest} \citep{maxentcppCompTest2025}. Its
contribution is not modernization in general, but optimizer-level
fidelity combined with integration into modern R/\texttt{terra}
workflows.

\subsection{Ecosystem comparison}\label{ecosystem-comparison}

The R ecosystem contains several packages that provide access to MaxEnt
models, each employing a distinct integration strategy:

{\def\LTcaptype{none} % do not increment counter
\begin{longtable}[]{@{}
  >{\raggedright\arraybackslash}p{(\linewidth - 4\tabcolsep) * \real{0.2250}}
  >{\raggedright\arraybackslash}p{(\linewidth - 4\tabcolsep) * \real{0.4750}}
  >{\raggedright\arraybackslash}p{(\linewidth - 4\tabcolsep) * \real{0.3000}}@{}}
\toprule\noalign{}
\begin{minipage}[b]{\linewidth}\raggedright
Package
\end{minipage} & \begin{minipage}[b]{\linewidth}\raggedright
MaxEnt integration
\end{minipage} & \begin{minipage}[b]{\linewidth}\raggedright
Dependency
\end{minipage} \\
\midrule\noalign{}
\endhead
\bottomrule\noalign{}
\endlastfoot
\textbf{dismo} \citep{Hijmans2023} & rJava bridge to \texttt{maxent.jar}
& Java JDK, rJava \\
\textbf{kuenm} \citep{Cobos2019} &
\texttt{system2("java\ -jar\ maxent.jar")} & Java JDK \\
\textbf{kuenm2} (Cobos et al.) & \texttt{glmnet} via forked
\texttt{maxnet} code & glmnet \\
\textbf{wallace} \citep{Kass2018wallace} & Delegates to ENMeval \(\to\)
maxnet or dismo & ENMeval \\
\textbf{ENMTools} \citep{Warren2010} & \texttt{dismo::maxent()} directly
& dismo, rJava \\
\textbf{biomod2} \citep{Thuiller2009} & Java (\texttt{system2}) or
\texttt{maxnet} & Optional Java or maxnet \\
\textbf{rmaxent} \citep{Baumgartner2017} & Java-free projection of
fitted Maxent models, \texttt{.lambdas} parsing, MESS & None (projection
only) \\
\textbf{MIAmaxent} \citep{Vollering2019} & Maximum entropy via subset
selection (not lasso); decoupled transformation/fitting/selection &
glmnet \\
\textbf{SDMtune} \citep{Vignali2020} & Trains/tunes SDMs via
\texttt{dismo::maxent()} or \texttt{maxnet}; hyperparameter tuning,
variable selection & dismo or maxnet \\
\textbf{maxentcpp} & Native C++17 via Rcpp & None (self-contained) \\
\end{longtable}
}

These packages can be grouped by their relationship to the MaxEnt
algorithm:

\begin{itemize}
\tightlist
\item
  \textbf{Black-box wrappers} (dismo, kuenm, ENMTools, biomod2 MAXENT
  mode): invoke the Java binary and read output files. The optimizer is
  inaccessible.
\item
  \textbf{Approximate reimplementations} (maxnet, kuenm2, biomod2 MAXNET
  mode): use \texttt{glmnet} coordinate descent on the logistic
  regression formulation. The statistical model is equivalent but the
  optimization path differs, which can produce numerically different
  results in some settings.
\item
  \textbf{Faithful reimplementation} (maxentcpp): ports the original
  Sequential optimizer to C++ and proves per-iteration numerical
  equivalence.
\end{itemize}

The distinguishing contribution of \texttt{maxentcpp} is that it is the
only package offering a compiled native reimplementation of the actual
MaxEnt \texttt{density.Sequential} optimizer --- preserving both the
statistical model and the optimization algorithm while eliminating the
Java dependency.

\subsection{Empirical comparison: maxentcpp vs
maxnet}\label{empirical-comparison-maxentcpp-vs-maxnet}

To quantify the practical differences between \texttt{maxentcpp} and
\texttt{maxnet}, we compared both packages on a virtual species
\citep{Leroy2016} with known true habitat suitability, constructed from
the environmental layers bundled with \texttt{maxentcpp} (Annual Mean
Temperature and Annual Precipitation over Central America, 2,371 non-NA
cells). The virtual species has a Gaussian niche centered at 20 °C and
1,500 mm, and 100 presence records were sampled proportionally to the
true suitability surface (see the companion vignette
\href{https://alrobles.github.io/maxentcpp/articles/virtual_species_comparison.html}{Virtual
Species Comparison} for full reproducible code).

Both packages were fitted with linear, quadratic, and hinge features. In
this illustrative example, the two packages produced highly similar rank
and spatial-overlap patterns:

{\def\LTcaptype{none} % do not increment counter
\begin{longtable}[]{@{}lr@{}}
\toprule\noalign{}
Metric & Value \\
\midrule\noalign{}
\endhead
\bottomrule\noalign{}
\endlastfoot
Schoener's \(D\) \citep{Schoener1968, Warren2008} & 0.93 \\
Spearman \(\rho\) (rank correlation) & 0.95 \\
Pearson \(r\) & 0.97 \\
Mean \(\lvert\Delta\text{cloglog}\rvert\) & 0.053 \\
Max \(\lvert\Delta\text{cloglog}\rvert\) & 0.43 \\
\end{longtable}
}

These results suggest that \texttt{maxentcpp} and \texttt{maxnet} can
produce highly similar predictions in a simple continuous-predictor
setting (Schoener's \(D > 0.9\) is conventionally interpreted as high
niche overlap), while also showing non-negligible local differences in
the tails of the suitability distribution, where the two optimization
algorithms (\texttt{maxentcpp}: sequential coordinate ascent;
\texttt{maxnet}: elastic-net coordinate descent) regularize differently.
A broader simulation study across multiple species, sample sizes, and
predictor correlation structures would be needed to establish general
equivalence.

The key practical scenarios where a user must choose \texttt{maxentcpp}
over \texttt{maxnet} are: (1) when exact reproduction of Java Maxent
results is required (e.g., replicating published analyses), (2) when
lambda file interoperability with Java Maxent is needed, and (3) when
streaming raster projection onto grids larger than available RAM is
required.

\subsection{Performance benchmarks}\label{performance-benchmarks}

Training time for the bundled \emph{Abeillia abeillei} dataset (73
presences, 2,371 background, linear + quadratic + hinge features):

{\def\LTcaptype{none} % do not increment counter
\begin{longtable}[]{@{}lr@{}}
\toprule\noalign{}
Package & Median time per fit \\
\midrule\noalign{}
\endhead
\bottomrule\noalign{}
\endlastfoot
\texttt{maxentcpp} & \textasciitilde820 ms \\
\texttt{maxnet} & \textasciitilde530 ms \\
\end{longtable}
}

\texttt{maxnet} is faster for small datasets because \texttt{glmnet}'s
coordinate descent is highly optimized for dense feature matrices.
\texttt{maxentcpp}'s streaming evaluation avoids materializing the full
prediction matrix, which is expected to reduce peak memory during
projection; future benchmarks should quantify this advantage on rasters
larger than available RAM.

To exercise the 1.0.0 feature set, we benchmarked the new diagnostics on
a synthetic grid of 2,371 cells and 100 presence records with two
continuous predictors plus one five-level categorical variable (linear +
quadratic + hinge features; \texttt{max\_iter\ =\ 500}):

{\def\LTcaptype{none} % do not increment counter
\begin{longtable}[]{@{}lr@{}}
\toprule\noalign{}
1.0.0 feature & Median wall time \\
\midrule\noalign{}
\endhead
\bottomrule\noalign{}
\endlastfoot
Single fit (continuous only) & \textasciitilde1 ms \\
Fit with categorical predictor & \textasciitilde94 ms \\
Jackknife (3 variables; 6 fits) & \textasciitilde75 s \\
Cross-validation (\(k = 5\); 5 fits) & \textasciitilde84 s \\
Replicate runs (bootstrap, \(n = 5\); 5 fits) & \textasciitilde108 s \\
\end{longtable}
}

Each 1.0.0 diagnostic is a small multiple of single fits, as expected:
jackknife runs one fit per variable per leave-out scheme, and
cross-validation and replicate runs each fit \(k\) or \(n\) models.
These wall times reflect the full per-fit cost on a mid-size grid and
are linear in the number of constituent fits; on the smaller bundled
\emph{Abeillia abeillei} dataset they are correspondingly faster. Full
reproducible code is provided in the package vignettes.

\section{Software design}\label{software-design}

\subsection{Architecture}\label{architecture}

\texttt{maxentcpp} was built as a new package rather than contributing
to existing projects for three reasons. First, a faithful port of the
original Java optimizer (sequential coordinate ascent using
\(\ell_1\)-penalized Newton steps) requires a C++ engine that cannot be
expressed through \texttt{glmnet}'s API; contributing this to
\texttt{maxnet} would change its core algorithm. Second,
\texttt{dismo}'s architecture is tightly coupled to \texttt{rJava} and
\texttt{raster}; the native C++ strategy eliminates this dependency
chain entirely. Third, the header-based C++ library design of
\texttt{maxentcpp} enables reuse in other packages or standalone
applications beyond R --- something not achievable through modifications
to existing R-only packages.

\texttt{maxentcpp} is organized in three layers (C++ core, Rcpp bridge,
R interface), with the C++ core further split into algorithmic and
infrastructure sublayers:

{\def\LTcaptype{none} % do not increment counter
\begin{longtable}[]{@{}
  >{\raggedright\arraybackslash}p{(\linewidth - 4\tabcolsep) * \real{0.2333}}
  >{\raggedright\arraybackslash}p{(\linewidth - 4\tabcolsep) * \real{0.5333}}
  >{\raggedleft\arraybackslash}p{(\linewidth - 4\tabcolsep) * \real{0.2333}}@{}}
\toprule\noalign{}
\begin{minipage}[b]{\linewidth}\raggedright
Layer
\end{minipage} & \begin{minipage}[b]{\linewidth}\raggedright
Key components
\end{minipage} & \begin{minipage}[b]{\linewidth}\raggedleft
Lines
\end{minipage} \\
\midrule\noalign{}
\endhead
\bottomrule\noalign{}
\endlastfoot
C++ core (algorithmic) & \texttt{Sequential}, \texttt{FeaturedSpace},
six feature types & \textasciitilde4,600 \\
C++ core (I/O, diagnostics) & \texttt{BackgroundProvider}, grid I/O,
CSV, MESS, response curves & \textasciitilde2,300 \\
Rcpp bridge & \texttt{rcpp\_*.cpp} external-pointer bindings
\citep{Eddelbuettel2013} & \textasciitilde2,300 \\
R interface & \texttt{maxent\_run()}, feature generators, projection,
evaluation, diagnostics, replicates, jackknife & \textasciitilde5,000 \\
\end{longtable}
}

The C++ core uses Eigen \citep{Guennebaud2010} for dense linear algebra.
Of the \textasciitilde6,900 C++ lines, approximately 4,600 implement the
core algorithmic logic (optimizer, features, density) while the
remainder handles I/O and diagnostics. The high-level
\texttt{maxent\_run()} function provides a one-call workflow mirroring
the Java GUI experience, while lower-level functions
(\texttt{maxent\_generate\_features()},
\texttt{maxent\_featured\_space()}, \texttt{maxent\_fit()},
\texttt{maxent\_project\_cloglog()}) offer full control over each
modeling step.

\subsection{Optimization algorithm}\label{optimization-algorithm}

The \texttt{Sequential} optimizer is a direct port of
\texttt{density.Sequential} in Java Maxent 3.4.4, as represented by the
public \texttt{mrmaxent/Maxent} repository and AMNH release notes (where
3.4.4 is described as a minor bug-fix release). It minimizes the
\(\ell_1\)-regularized negative log-likelihood:

\[
\mathcal{L}(\boldsymbol{\lambda}) = -\frac{1}{m}\sum_{i=1}^{m}
\sum_{j=1}^{J} \lambda_j f_j(\mathbf{x}_i) + \log Z(\boldsymbol{\lambda})
+ \sum_{j=1}^{J} \beta_j |\lambda_j|,
\]

where \(m\) is the number of presence records, \(Z(\boldsymbol{\lambda})
= \sum_{k=1}^{N} \exp\!\bigl(\sum_j \lambda_j f_j(\mathbf{x}_k)\bigr)\)
is the partition function summed over \(N\) background points, and
\(\beta_j = \hat\sigma_j / \sqrt{m}\) is the per-feature regularization
parameter with \(\hat\sigma_j\) being the standard deviation of feature
\(j\) over the presence points (floored at 0.001).

At each iteration, the optimizer selects the feature \(j^*\) with the
most negative \(\Delta\mathcal{L}\) bound (the \texttt{deltaLossBound}
function; block-coordinate ascent \citep{Tseng2001}), then computes a
Newton step:

\[
\alpha_j = -\frac{\partial \mathcal{L}/\partial \lambda_j}
{\mathbf{u}_j^\top \mathbf{H} \mathbf{u}_j},
\]

where \(\mathbf{u}_j\) is the \(j\)-th coordinate direction and
\(\mathbf{u}_j^\top \mathbf{H} \mathbf{u}_j = \mathrm{Var}_q[f_j]\) is
the variance of feature \(j\) under the current distribution \(q\). The
derivative includes the \(\ell_1\) subgradient:
\(\partial \mathcal{L}/\partial \lambda_j = \mathbb{E}_q[f_j] -
\mathbb{E}_{\hat p}[f_j] + \beta_j \,\mathrm{sign}(\lambda_j)\). The
step is damped during early iterations (\(\alpha \leftarrow \alpha/50\)
for iterations \(<10\), \(\alpha/10\) for \(<20\), \(\alpha/3\) for
\(<50\)) and falls back to a line search (\texttt{searchAlpha}) if the
Newton step does not decrease loss. Every 10 iterations, a parallel
update applies coordinate steps to all features simultaneously.

Convergence is tested every 20 iterations: training stops when the
relative loss improvement falls below \(10^{-5}\) or 500 iterations are
reached. All constants in this section are taken directly from Java
Maxent 3.4.4 source files, and every numerical constant and control-flow
branch is covered by a source-mapping test in
\texttt{maxentcppCompTest}.

\subsection{Streaming raster
evaluation}\label{streaming-raster-evaluation}

\texttt{maxentcpp} reads \texttt{terra::SpatRaster} objects
block-by-block through a \texttt{BackgroundProvider} abstraction. Each
tile is scored by the C++ engine and discarded before the next tile is
loaded, allowing projection onto raster stacks larger than available
RAM. The partition function \(Z\) is accumulated across tiles by summing
unnormalized densities \(\exp(\sum_j \lambda_j f_j(\mathbf{x}_k))\) for
each cell \(k\) in the tile, then normalizing once after all tiles are
processed. This produces results equivalent to single-pass evaluation
because the mathematical sum is commutative and each cell is scored
independently; floating-point summation order may introduce differences
at the least-significant bit level.

\subsection{Numerical fidelity}\label{numerical-fidelity}

Each C++ class is a direct translation of its Java counterpart,
preserving the same control flow, regularization logic, and numerical
constants:

{\def\LTcaptype{none} % do not increment counter
\begin{longtable}[]{@{}ll@{}}
\toprule\noalign{}
\texttt{maxentcpp} (C++) & Java Maxent original \\
\midrule\noalign{}
\endhead
\bottomrule\noalign{}
\endlastfoot
\texttt{LinearFeature} & \texttt{density.LinearFeature} \\
\texttt{QuadraticFeature} & \texttt{density.SquareFeature} \\
\texttt{ProductFeature} & \texttt{density.ProductFeature} \\
\texttt{HingeFeature} & \texttt{density.HingeFeature} \\
\texttt{ThresholdFeature} & \texttt{density.ThresholdFeature} \\
\texttt{BinaryFeature} & \texttt{density.BinaryFeature} \\
\texttt{FeaturedSpace} & \texttt{density.FeaturedSpace} \\
\texttt{Sequential::run()} & \texttt{density.Sequential.run()} \\
\end{longtable}
}

This design choice prioritizes backward compatibility: users migrating
from Java Maxent should obtain numerically equivalent results for
implemented features under matched inputs, defaults, and deterministic
settings. The fidelity contract is enforced by a dedicated companion
package, \texttt{maxentcppCompTest} \citep{maxentcppCompTest2025}, which
runs the C++ optimizer and the original Java \texttt{density.Sequential}
on shared test fixtures and compares per-iteration trajectories of loss,
entropy, and \(\lambda\) vectors.

On controlled test fixtures (identical input data, same floating-point
representation), agreement reaches \(< 10^{-14}\) relative error for
symmetric fixtures and \(< 10^{-6}\) on \(\|\Delta\lambda\|_\infty\) for
asymmetric fixtures at every iteration checkpoint. \textbf{These bounds
hold under the specific conditions of the test fixtures} (same compiler
family, IEEE 754 double precision, deterministic iteration order).
Floating-point ordering differences introduced by alternative BLAS
backends or aggressive compiler optimizations (e.g.,
\texttt{-ffast-math}) could widen the gap, though the sequential nature
of the coordinate-ascent algorithm limits sensitivity to
parallelism-induced reordering. Cross-platform CI (Ubuntu, macOS,
Windows with GCC, Clang, and MSVC) confirms that the test fixtures pass
on all three compiler families.

\textbf{Lambda file compatibility.} Trained models are serialized in the
same \texttt{.lambdas} text format used by Java Maxent 3.4.4. Lambda
files produced by either implementation can be loaded by the other,
ensuring interoperability with existing workflows and archived models.

\subsection{Feature completeness}\label{feature-completeness}

The following table summarizes the current scope of \texttt{maxentcpp}
relative to Java Maxent 3.4.4 and \texttt{maxnet}:

Legend: ✓ = implemented and tested; ◐ = partial or experimental; --- =
not implemented.

{\def\LTcaptype{none} % do not increment counter
\begin{longtable}[]{@{}lccc@{}}
\toprule\noalign{}
Feature & \texttt{maxentcpp} & Java Maxent 3.4.4 & \texttt{maxnet}
0.1.4 \\
\midrule\noalign{}
\endhead
\bottomrule\noalign{}
\endlastfoot
Linear features & ✓ & ✓ & ✓ \\
Quadratic features & ✓ & ✓ & ✓ \\
Product features & ✓ & ✓ & ✓ \\
Hinge features & ✓ & ✓ & ✓ \\
Threshold features & ✓ & ✓ & ✓ \\
Raw/logistic/cloglog output & ✓ & ✓ & ✓ \\
AUC evaluation & ✓ & ✓ & via ENMeval \\
Permutation importance & ✓ & ✓ & via ENMeval \\
Response curves & ✓ & ✓ & ◐ \\
MESS maps & ✓ & ✓ & --- \\
Clamping / fade-by-clamping \citep{Phillips2008} & ✓ & ✓ & --- \\
Lambda file I/O & ✓ & ✓ & --- \\
Streaming raster projection & ✓ & --- & --- \\
Bias grids & ✓ & ✓ & via \texttt{offset} \\
Categorical variables & ✓ & ✓ & ✓ \\
Replicate runs / cross-validation & ✓ & ✓ & via ENMeval \\
Jackknife variable selection & ✓ & ✓ & --- \\
Missing data handling & ✓ & ✓ & ✓ \\
SWD-to-raster projection & --- & ✓ & --- \\
\end{longtable}
}

Categorical variables, replicate runs, jackknife, missing-data handling,
and background bias weighting are now implemented and tested.
SWD-to-raster projection remains a workflow convenience not yet exposed
through the public R API; users can still project models by loading
environmental grids via \texttt{terra} and the package's grid conversion
functions.

\subsection{Development provenance}\label{development-provenance}

The translation followed a phased approach across four publicly
accessible repositories:

\begin{enumerate}
\def\labelenumi{\arabic{enumi}.}
\tightlist
\item
  \texttt{alrobles/Maxent} --- a fork of the original
  \texttt{mrmaxent/Maxent} \citep{Phillips2017} Java source code, where
  the architecture was analyzed and each Java class was mapped to a C++
  target.
\item
  \texttt{alrobles/maxentcpp-devel} --- the development repository where
  the C++ port, Rcpp bindings, R interface, tests, and documentation
  were iteratively built and refined.
\item
  \texttt{alrobles/maxentcppCompTest} --- the cross-language fidelity
  test suite containing Java oracle runners, golden trajectory files,
  and automated comparison scripts.
\item
  \texttt{alrobles/maxentcpp} --- the release repository with the clean,
  CRAN-ready package.
\end{enumerate}

\section{Limitations and inherited
assumptions}\label{limitations-and-inherited-assumptions}

The Maxent 3.4.4 algorithm that \texttt{maxentcpp} faithfully reproduces
has well-documented limitations that users should be aware of:

\begin{itemize}
\tightlist
\item
  \textbf{Feature explosion.} The number of features (particularly hinge
  and threshold) grows with the number of environmental variables, which
  can cause identifiability issues in high-dimensional settings
  \citep{Elith2011, Merow2013}.
\item
  \textbf{Sensitivity to background sampling.} Maxent's output is
  influenced by the choice and size of the background sample, which
  affects the estimated density and can bias predictions in undersampled
  regions \citep{Phillips2009}.
\item
  \textbf{\(\ell_1\) regularization limitations.} The sequential
  coordinate ascent with \(\ell_1\) penalties produces sparse models but
  does not guarantee selection of the ``correct'' features under high
  correlation among predictors \citep{Hastie2015}.
\item
  \textbf{No principled uncertainty quantification.} Unlike Bayesian
  approaches or bootstrap-based methods, the Maxent optimizer produces
  point estimates without confidence intervals.
\end{itemize}

A principled alternative would be to reformulate the Maxent objective
using modern convex optimization tooling (e.g., proximal gradient
methods, ADMM) or Bayesian inference. However, such reformulations would
produce different results than the widely used Java implementation,
breaking comparability with published analyses. The design of
\texttt{maxentcpp} explicitly aims to preserve this comparability.

\section{CRAN compliance and software
quality}\label{cran-compliance-and-software-quality}

Achieving CRAN acceptance requires more than passing
\texttt{R\ CMD\ check}: packages must meet strict software-quality
standards that safeguard user environments and ensure reproducibility
across platforms \citep{CRANPolicy2024, WRE2024}. During the
pre-submission review, \texttt{maxentcpp} was brought to compliance on
four fronts. First, all 36 example blocks were migrated from
\texttt{\textbackslash{}dontrun\{\}} to
\texttt{\textbackslash{}donttest\{\}} \citep{CRANcookbook2024} and
rewritten to be self-contained with synthetic data, so every documented
example is live, testable code that CRAN's infrastructure can execute
with \texttt{-\/-run-donttest}. Second, functions that touch global
graphics state now restore it through an internal
\texttt{.maxent\_safe\_png()} helper built on \texttt{on.exit()},
guaranteeing cleanup even on error \citep{WRE2024}. Third, all
\texttt{terra}-dependent examples are guarded with
\texttt{requireNamespace("terra",\ quietly\ =\ TRUE)} so they skip
gracefully when the optional dependency is unavailable. Fourth, GitHub
Actions runs \texttt{R\ CMD\ check\ -\/-as-cran} on Ubuntu (R release
and R-devel), macOS, and Windows, plus a CRAN-preflight job with
\texttt{\_R\_CHECK\_FORCE\_SUGGESTS\_=false} that simulates CRAN's
minimal-dependency environment.

\section{Research impact statement}\label{research-impact-statement}

\texttt{maxentcpp} provides a numerically faithful implementation of the
core continuous-predictor Maxent fitting and projection workflow for the
implemented feature classes and output transformations, including
categorical predictors, replicate runs and cross-validation, jackknife
diagnostics, and missing-data handling. SWD-to-raster projection remains
a convenience workflow not yet exposed through the public R API.

The package is distributed under the MIT license (compatible with
Eigen's MPL2 and RcppEigen's GPL-2 via the ``or later'' clause), with
full source code, a pkgdown documentation website at
\url{https://alrobles.github.io/maxentcpp/}, and four vignettes covering
a getting-started walkthrough, the mathematical foundations of Maxent
features, a comparative motivation document, and a virtual species
validation study. Continuous integration via GitHub Actions runs
\texttt{R\ CMD\ check} on Ubuntu, macOS, and Windows across R release
and R-devel. The test suite comprises over 20 test files
(\textasciitilde4,800 lines) covering feature generation, optimization
convergence, spatial projection, Java numerical equivalency, model
diagnostics, and end-to-end workflows.

By providing a dependency-free, memory-efficient, and numerically
faithful Maxent implementation, \texttt{maxentcpp} may lower the barrier
for large-scale biodiversity assessments, climate-change projections,
and conservation prioritization studies that rely on presence-only
species distribution models, bringing the package to near-parity with
Java Maxent 3.4.4 for the implemented feature classes while remaining
installable in Java-free environments.

\section{AI usage disclosure}\label{ai-usage-disclosure}

Generative AI tools were used during the development of
\texttt{maxentcpp} and the preparation of this manuscript, as detailed
below.

\textbf{Tools used.}

\begin{itemize}
\tightlist
\item
  \textbf{GitHub Copilot} (GPT-4-based, 2025 version): code scaffolding
  and boilerplate generation during the initial porting phases in
  \texttt{alrobles/Maxent} and \texttt{alrobles/maxentcpp-devel}.
\item
  \textbf{Devin} (Cognition AI, 2025--2026): incremental feature
  implementation, test writing, documentation generation, CRAN
  compliance fixes, CI configuration, and manuscript drafting.
\item
  \textbf{Claude} (Anthropic, 2025--2026): complex refactoring,
  cross-cutting documentation, and code review.
\end{itemize}

\textbf{Nature and scope of assistance.} AI tools contributed to: C++
code generation and refactoring from Java source, Rcpp binding
scaffolding, R wrapper functions, \texttt{testthat} test scaffolding and
expansion, \texttt{roxygen2} documentation, vignette drafting, pkgdown
site configuration, CI workflow setup, CRAN compliance review, and
manuscript drafting. The core algorithmic C++ classes
(\texttt{Sequential}, \texttt{FeaturedSpace}, feature types) were
translated from Java with AI assistance but under direct human
supervision: each class was mapped line-by-line from the corresponding
Java source, reviewed for correctness against the Java reference, and
validated through the \texttt{maxentcppCompTest} trajectory comparison
framework. Approximately 60\% of the C++ algorithmic lines were
initially drafted by AI tools and subsequently reviewed and corrected by
the human author; the remaining 40\% were written directly by the
author, particularly the performance-critical optimizer inner loops and
numerical edge cases.

\textbf{Maintainability.} The human author (Á. L. Robles Fernández) can
independently explain and modify every algorithmically significant
class. As evidence: the \texttt{Sequential::run()} optimizer,
\texttt{good\_alpha()}, \texttt{delta\_loss\_bound()},
\texttt{newton\_step\_feature()}, and \texttt{do\_parallel\_update()}
methods are documented with line-by-line references to the corresponding
Java source (e.g., ``mirrors Sequential.java:294..342''), and the author
has independently debugged numerical discrepancies during the fidelity
testing process.

\textbf{Confirmation of review.} All AI-generated code, tests,
documentation, and manuscript text were reviewed, edited, and validated
by the human author, who made all core design decisions including the
algorithm translation strategy, API design, numerical fidelity
requirements, and the phased development architecture. The full
development history is publicly available in the four repositories
listed above.

\section{Acknowledgements}\label{acknowledgements}

The author thanks Steven J. Phillips, Robert P. Anderson, and Robert E.
Schapire for creating the original Maxent software and making the source
code publicly available under an MIT license. The \texttt{Rcpp},
\texttt{RcppEigen}, and \texttt{terra} package maintainers are
gratefully acknowledged for providing the infrastructure that makes
\texttt{maxentcpp} possible.

\renewcommand\refname{References}
\bibliography{paper.bib}

\end{document}
